\section{三行主计数定理的形式化验证}\label{sec:formalization}

\subsection{数学对象与定理的对应}
本版将审阅范围固定为 M1--M6；后续本地阶段另行审阅。
本节说明定理 \ref{thm:main} 的机器验证范围。
形式化沿第 \ref{sec:half} 节的半棋盘路线进行，
不依赖第 \ref{sec:blocks} 节的基本块分类。
代码中的棋盘格组成有限类型，砖是两个共边相邻格的无序集合，
铺法要求每格恰由一块砖覆盖。
四元群通过行反转、列反转及其复合在这些铺法上作用。
记号 \texttt{a n} 定义为 $3\times2n$ 铺法的该作用轨道数；
\texttt{orbitSeries} 的第 $n$ 项系数为 $(a(n):\mathbb Q)$，
并非由目标递推或预设有理式定义的辅助数列。

设
\[
 B(x)=(1-x)(1-4x+x^2)(1-4x^2+x^4),\qquad
 N(x)=1-3x-4x^2+10x^3-4x^5+x^6.
\]
形式化首先证明 $B(x)A(x)=N(x)$，再由 $B(0)=1$ 证明 $B$ 是
$\mathbb Q[[x]]$ 中的单位，从而得到 $A=NB^{-1}$。
随后比较 $n+7$ 次系数，在整数中证明对每个 $n\ge0$ 有
\[
\begin{split}
 a(n+7)={}&5a(n+6)-a(n+5)-19a(n+4)+19a(n+3)\\
          &+a(n+2)-5a(n+1)+a(n).
\end{split}
\]
比较 $0,\ldots,6$ 次系数得到定理 \ref{thm:main} 的全部初值，
并进一步验证递推起点 $a(7)=2037$。
含负项的等式在整数或有理数中解释，另给出自然数中的移项版本，
因而没有使用自然数截断减法代替整数减法。

\begin{table}[htbp]
\centering
\caption{正文结论与主要 Lean 声明的对应。}\label{tab:lean}
\small
\begin{tabular}{ll}
\toprule
正文结论 & \texttt{Domino} 命名空间中的声明\\
\midrule
四类固定点生成函数 & \texttt{fixed\_series\_identities}\\
主分母恒等式 $BA=N$ & \texttt{orbit\_series\_identity}\\
式 \eqref{eq:main} 的分式版本 & \texttt{orbit\_series\_rational}\\
原题的符号排列 & \texttt{orbit\_series\_paper}\\
全部指标七阶递推 & \texttt{orbit\_recurrence\_seven}\\
七项初值 & \texttt{orbit\_initial\_values}\\
递推起点 $a(7)=2037$ & \texttt{orbit\_recurrence\_first}\\
主定理汇总 & \texttt{three\_row\_main\_theorem}\\
\bottomrule
\end{tabular}
\end{table}

证明链依次建立真实铺法与状态路径的双射、两种中线补全及限制的互逆关系、
三行上下反射固定铺法的唯一性、任意步数的状态计数递推、
真实轨道的 Burnside 公式，以及系数级数恒等式。
几何双射、全参数归纳和系数比较均由证明项承担；
有限状态判定只用于局部规则等有限命题。

\subsection{验证环境与覆盖边界}
本项目固定 Lean 4.35.0 和 mathlib 提交
\texttt{56a005855ea17e0b0958f73ec801ae140b0a1679}。
2026 年 10 月 25 日完成的 M1--M6 验收包含完整项目构建、
32 个数学模块中 413 项显式声明的累计公理审计，
以及项目源码在无自身旧编译产物目录中的重新构建。
审计只出现 \texttt{propext}、\texttt{Classical.choice} 和 \texttt{Quot.sound}
三个标准逻辑公理，未使用占位证明或自行加入的公理。
413 项声明包含辅助定义和实例，不等于论文定理的数量。
重新构建复用了已核对版本的第三方依赖缓存，
不表示对全部 mathlib 进行了无缓存编译。

上述验证覆盖定理 \ref{thm:main} 及其三行半棋盘计数依赖。
精确稳定子分类、渐近结论、一般奇数高度公式、五行证书和基本块路线
具有本文所给的书面证明，但尚未纳入这一形式化验收。
代码明确使用四种坐标对称，没有另行形式化欧氏平面全部等距变换的分类。
零宽采用抽象四元群对唯一空铺法的平凡作用。
机器验证不判定文献首次性；原题与形式化定义的对应仍应接受数学审阅。


\subsection{语义审阅入口}
仓库的 \path{review/m1-m6/} 保存逐项对应表、依赖说明、源码摘要、
可执行的声明检查以及待外部填写的审阅表。
审阅时应先读定义，再看双射与计数桥接，最后检查主定理类型。
尤其需要区分：Lean 中 $\alpha,\beta$ 先按递推定义，
再由 \texttt{profile\_even\_odd} 对任意宽度证明它们等于几何基数。
只有递推定义本身不能证明它们计数论文对象。
代码执行成功和人工确认题意一致是两种不同证据，审阅表分别记录。
